This paper presented a case study on hardening an existing open-source microservices application. By establishing a stabilized Docker environment and a GitHub Actions CI pipeline, we created a reproducible foundation for quality assurance.

Our empirical results show that while standard unit testing provided a starting point, it left significant gaps in branch coverage (only 14.4\% covered). The introduction of Mutation Testing was critical in exposing these gaps, driving improvements in the domain layer where semantic logic resides. However, our attempt to introduce formal verification via JML highlighted a significant gap in the ecosystem: modern frameworks like Spring Boot and Lombok are currently resistant to standard formal verification tools due to their reliance on runtime injection and bytecode generation.

We conclude that for existing "modern" Java applications, a combination of \textbf{Containerization}, \textbf{Strict Static Analysis}, and \textbf{Targeted Mutation Testing} offers the highest return on investment for dependability, while formal verification remains a specialized technique best suited for core algorithmic modules rather than full-stack web architectures.